\chapter{Research Methodology}

%% English draft. Uses \crosscheck{...} (defined in main.tex) to flag content that a
%% human must verify before submission. All \cite{...} keys resolve to references.bib.
%% Figures are TikZ (tikz/pgfplots are loaded by thesisdtetiugm.cls).

\section{Research Overview}

This research follows \emph{design science research} (DSR), the paradigm for work whose
contribution is a constructed artifact rather than a tested hypothesis. Where behavioural
research seeks to explain or predict, design science seeks to extend capability by building
and evaluating an artifact, and knowledge of the problem is obtained \emph{through} that
construction \cite{hevner2004}. The specific process model followed is the design science
research methodology of Peffers \textit{et al.} \cite{peffers2007}, whose six activities are
problem identification and motivation, definition of objectives for a solution, design and
development, demonstration, evaluation, and communication.

The artifact here is a FHIR-native, knowledge-graph based clinical recommendation system
(CRS) for chronic coronary disease (CCD); its evaluation is the concordance, agreement,
safety, and faithfulness assessment of Section~\ref{sec:eval-method}. The six DSR activities map onto this thesis as follows. Problem identification and
objectives are Chapter~I. Design and development are the four stages below. Demonstration is
the grounded briefing produced for individual patients. Evaluation is Chapter~IV, and
communication is the thesis itself. The
work is organised into four stages:

\begin{enumerate}
    \item \textbf{Data generation.} Synthetic patient records are generated with Synthea
    and exported as HL7 FHIR R4 bundles.
    \item \textbf{Knowledge-graph construction.} The FHIR bundles are loaded into a Neo4j
    knowledge graph, which is then fused with the 2023 AHA/ACC CCD guideline and a
    guideline-seeded biomedical-literature corpus through a UMLS-based concept layer.
    \item \textbf{Recommendation-system implementation.} A query layer matches each patient to
    applicable guideline recommendations, gates observation-dependent recommendations on
    measured values, detects contraindication-safety flags, ranks literature-supported
    options, and produces a grounded natural-language briefing.
    \item \textbf{Evaluation.} The system is evaluated for graph construction quality,
    guideline concordance (with an observation-gate ablation), agreement with a manually
    labelled gold standard, safety detection, and briefing faithfulness.
\end{enumerate}

Figure~\ref{fig:research-flow} summarises the flow and the principal artifact produced by
each stage. Each stage is elaborated in its own section below: stage~1 in
Section~\ref{sec:data-source}, stage~2 in Section~\ref{sec:kg-construction}, stage~3 in
Section~\ref{sec:dss-method}, and stage~4 in Section~\ref{sec:eval-method}.

\begin{figure}[htbp]
\centering
\begin{tikzpicture}[
    node distance = 7mm and 12mm,
    stage/.style = {rectangle, rounded corners=2pt, draw, thick,
                    text width=62mm, align=center, minimum height=11mm,
                    font=\small},
    art/.style   = {rectangle, draw, dashed, text width=32mm, align=center,
                    minimum height=8mm, font=\scriptsize, fill=black!5},
    ar/.style    = {-latex, thick}
]
\node[stage] (s1)
    {\textbf{1. Data generation}\\[1pt]
     \scriptsize Synthea $\rightarrow$ HL7 FHIR R4 bundles};
\node[stage, below=of s1] (s2)
    {\textbf{2. Knowledge-graph construction}\\[1pt]
     \scriptsize FHIR $\rightarrow$ Neo4j; UMLS concept bridge;
     guideline and literature layers};
\node[stage, below=of s2] (s3)
    {\textbf{3. Recommendation-system implementation}\\[1pt]
     \scriptsize applicable set, observation gates, safety flags,
     literature ranking, grounded briefing};
\node[stage, below=of s3] (s4)
    {\textbf{4. Evaluation}\\[1pt]
     \scriptsize graph structure, guideline concordance and ablation,
     gold standard, safety, faithfulness};

\node[art, right=of s1] (a1) {FHIR bundles \\ and CSV export};
\node[art, right=of s2] (a2) {Fused knowledge graph};
\node[art, right=of s3] (a3) {Recommendations, safety flags, briefing};
\node[art, right=of s4] (a4) {Evaluation metrics};

\draw[ar] (s1) -- (s2);
\draw[ar] (s2) -- (s3);
\draw[ar] (s3) -- (s4);
\draw[ar, dashed] (s1) -- (a1);
\draw[ar, dashed] (s2) -- (a2);
\draw[ar, dashed] (s3) -- (a3);
\draw[ar, dashed] (s4) -- (a4);
\end{tikzpicture}
\caption{Research flow: the four stages of the study and the principal artifact
produced by each.}
\label{fig:research-flow}
\end{figure}

\section{Tools and Materials}
\label{sec:tools}

The software and materials used in this research are listed below. Versions and minimum
specifications are stated so that another researcher can reproduce the build on
comparable hardware.

\begin{itemize}
    \item \textbf{Synthea} synthetic patient generator \cite{synthea}, build
    \texttt{aa0772f}, running on Java~26.0.1, producing HL7 FHIR~R4 bundles
    \cite{fhir_r4}. The generation parameters are recorded in
    Section~\ref{sec:cohort}.
    \item \textbf{Neo4j} graph database, Enterprise edition~2026.05.0, as the
    knowledge-graph store \cite{neo4j}, queried with Cypher \cite{cypher}. The instance runs
    under Neo4j Desktop~2.2.1 and is configured with a JVM heap of 4~GB initial and 8~GB
    maximum and a 4~GB page cache; the graph is held in a database named
    \texttt{sihdfhir}.
    \item \textbf{Python} for the loader, terminology bridge, recommendation queries, and
    evaluation scripts, using the official Neo4j driver.
    Two interpreters are used: Python~3.14.4 for the loader, recommendation-system and
    evaluation scripts (\texttt{neo4j}~6.2.0, \texttt{anthropic}~0.112.0,
    \texttt{biopython}~1.87), and Python~3.11.9 for the natural
    language stack, which \texttt{nmslib} pins below 3.12 (\texttt{scispacy}~0.6.2,
    \texttt{spacy}~3.7.5, \texttt{numpy}~1.26.4, \texttt{scikit-learn}~1.8.0,
    \texttt{pandas}~3.0.3).
    \item \textbf{scispaCy} clinical NLP models (\texttt{en\_core\_sci\_md}) with UMLS
    \cite{umls} and MeSH entity linkers \cite{scispacy}, used for literature entity
    extraction and terminology bridging.
    The model is \texttt{en\_core\_sci\_md}~0.5.4; the linkers are built over the UMLS
    2022AB subset (semantic categories 0, 1, 2 and 9) and the 2022 MeSH release. Three
    similarity thresholds are used, tightening with the precision each stage demands:
    0.80 for literature entity linking, 0.85 for clinical entity linking, and 0.75 for the
    SNOMED~CT-to-UMLS bridge, where concept names link more reliably than free text.
    \item \textbf{Large language model(s)} accessed via API, used only to phrase the
    retrieved recommendation facts into a briefing.
    The briefing reported here is generated with \texttt{claude-haiku-4-5}; a parallel
    implementation using \texttt{gpt-4o-mini} produces the same briefing from the same
    retrieved facts. The choice of model affects only the phrasing: retrieval is
    deterministic and completes before any model is called, so the facts a briefing may
    contain are fixed before generation begins.
    \item \textbf{Hardware.} A workstation capable of holding the graph working set in
    memory is required; the dominant constraint is the Neo4j page cache during the bulk
    load rather than CPU.
    The work reported here ran on a 13th-generation Intel Core i7-1365U at 1.8~GHz with
    32~GB of RAM under Windows~11 Pro. The database reserves 12~GB on its own: an
    8~GB maximum heap and a 4~GB page cache. 16~GB is therefore the practical floor for
    reproducing the build, and below that the page cache thrashes during the bulk load.
    32~GB leaves room for the load and the evaluation scripts to run together.
\end{itemize}

\section{Data Source and Study Cohort}
\label{sec:data-source}

\subsection{Data source}
The patient data are entirely \emph{synthetic}, generated by Synthea \cite{synthea} and
exported as HL7 FHIR R4 bundles \cite{fhir_r4}; no real patient data are used, which
minimises privacy and ethical risk. In the FHIR bundles, conditions and procedures are
coded in SNOMED~CT \cite{snomedct}, medications in RxNorm \cite{rxnorm}, and observations
(laboratory results and vital signs) in LOINC \cite{loinc}; ICD-10-CM/ICD-11 codes are not
present.

A FHIR \texttt{Bundle} carries no \texttt{fhirVersion} element---only a
\texttt{CapabilityStatement} does, and Synthea emits none---so the release was confirmed
from the exported instances rather than assumed. Three independent structural
discriminators agree on R4. First, coded concepts resolve to the
\texttt{terminology.hl7.org/CodeSystem/} namespace with no occurrence of the pre-R4
\texttt{hl7.org/fhir/v3/} or \texttt{/v2/} namespaces. Second, claim line items use
\texttt{productOrService}, the element R4 introduced in place of the STU3
\texttt{Claim.item.service}. Third, medication orders appear as
\texttt{MedicationRequest} rather than the DSTU2 \texttt{MedicationOrder}. The exporter's
output directory (\texttt{fhir/}, as opposed to \texttt{fhir\_stu3/} or
\texttt{fhir\_dstu2/}) is consistent with this. The bundles additionally carry US~Core
profile assertions, and US~Core is defined only against R4.
The same instance-level evidence dates the implementation guide. The profile set includes
\texttt{us-core-observation-screening-assessment}, introduced in US~Core~6.0.0 to supersede
the earlier Survey profile, which is itself absent; \texttt{us-core-condition-encounter-diagnosis}
is present while the pre-5.0.0 single \texttt{us-core-condition} profile is not. Synthea
offers four US~Core versions---3.1.1, 4.0.0, 5.0.1 and 6.1.0 \cite{synthea_fhir_wiki}---of
which only 6.1.0 is at or above 6.0.0. The export therefore used US~Core~6.1.0. No
\texttt{synthea.properties} was retained from the generating build, so this rests on the
exported instances, on the same footing as the release determination above.

\subsection{Cohort}
\label{sec:cohort}
The dataset was produced by two separate Synthea runs, both for the state of Massachusetts
with ten years of exported history, and combined into a single corpus
(Table~\ref{tab:generation}). The first run was restricted to the \texttt{stable\_ischemic\_heart\_disease} module so that
the corpus would contain a sufficient density of patients with the target disease. The
second ran the full module set (\texttt{*}) to provide a general background population.
Cohort selection therefore has to discriminate, rather than being trivially satisfied by
every record. The two runs yield
$2{,}009 + 3{,}481 = 5{,}490$ patients in total. Each run recorded its random seed, so the
corpus is exactly regenerable.

\begin{table}[h]
\centering
\caption{Synthea generation parameters for the two runs comprising the dataset.}
\label{tab:generation}
\begin{tabular}{lll}
\hline
\textbf{Parameter} & \textbf{Disease-enriched run} & \textbf{General run} \\
\hline
Modules            & \texttt{stable\_ischemic\_} & \texttt{*} (all) \\
                   & \texttt{heart\_disease} & \\
Patients           & 2{,}009 & 3{,}481 \\
Seed               & 1784186620874 & 1784186859476 \\
State              & Massachusetts & Massachusetts \\
Years of history   & 10 & 10 \\
Synthea build      & \texttt{aa0772f} & \texttt{aa0772f} \\
\hline
\multicolumn{3}{l}{\textbf{Total: 5{,}490 patients}} \\
\hline
\end{tabular}
\end{table}

\subsection{Why module selection does not determine cohort size}
\label{sec:module-yield}
It is tempting to read the 2{,}009 patients of the disease-enriched run as 2{,}009 patients
with chronic coronary disease. They are not. A Synthea module is a \emph{probabilistic
state machine}, not a label: enabling
\texttt{stable\_ischemic\_heart\_disease} means every patient in that run is
\emph{simulated through} the module, not that every patient reaches its diagnosis state.
Progression through the module is driven by age, sex, race, and accumulated risk factors such
as smoking status, diabetes, hypertension, and lipid levels, all of which are themselves
generated stochastically. Figure~\ref{fig:sihd-module} gives the module's full state machine;
the mechanism is worth stating precisely, because it explains both why the yield is partial
and why the diagnosed patients are middle-aged.

The module itself contains no age condition. Its \texttt{Chance\_of\_IHD} state is a one-month
delay whose transition to onset carries the \texttt{ihd\_risk} attribute as its probability,
defaulting to zero, and which otherwise loops back on itself. Every simulated month is an
independent draw, and a patient may loop for an entire lifetime without ever transitioning.
The age dependence enters through \texttt{ihd\_risk} itself, which Synthea computes as a
ten-year cardiovascular risk converted to a per-month hazard and multiplied by five. In the
build used here that risk comes from the ASCVD Pooled Cohort Equations rather than from
Framingham---the \texttt{useFramingham} flag is \texttt{false}---and
\texttt{ASCVD.ascvd10Year} returns $-1$ instead of a probability whenever the patient is
younger than 40 or older than 79. Outside that window no onset can occur. Inside it, age
enters continuously through $\ln(\mathrm{age})$ and its interaction terms, so the monthly
hazard climbs across the lifespan.

The corpus bears this out on both counts. Of the 876 recorded onsets the youngest
is at age 41 and the median at 60, with none below 40 and exactly one
above 79---the calculator's window reproduced exactly in the data. The age-40 floor is
therefore not a clinical judgement encoded in the module but an artefact of the risk
calculator's validated range, and is better read as a limitation of the synthetic cohort than
as a modelling choice. Yield tells the same story, and can be measured exactly rather than
bounded. The two runs share no patient identifier, and the disease-enriched run emits no
tobacco-status observation at all while every patient in the general run carries one, so the
presence of LOINC \texttt{72166-2} identifies a patient's run of origin unambiguously. On that
marker the disease-enriched run contributes 249 of the 876 and the general run
627---yields of 12.4\% ($249/2{,}009$) and 18.0\%
($627/3{,}481$). The disease-enriched run's own \texttt{conditions.csv} confirms the first
figure independently: it contains exactly 249 rows carrying \texttt{414545008}.

The run selected \emph{for} the disease therefore produced proportionally \emph{fewer} cases
than the run that merely enabled every module. Module selection did not simply fail to
determine cohort size; it moved it in the direction opposite to the one intended. The cause is
visible in the same export. Across its 2{,}009 patients the disease-enriched run records only
three distinct condition codes and 606 condition rows in total, with
no hypertension, diabetes or hyperlipidaemia anywhere, against 285
codes and 130{,}859 rows in the general run. Onset is driven by the ASCVD risk
score, and that score integrates precisely those comorbidities: enabling the disease module
alone starves its own risk model of the risk factors it depends on. Disease prevalence in a
Synthea corpus is thus an \emph{outcome} of the simulation, to be measured after generation,
not a parameter that can be set in advance.

One consequence carries into the evaluation, and Chapter~IV takes it up. The 249
disease-enriched patients in the cohort carry an ischemic-heart-disease diagnosis, an abnormal
cardiac imaging finding, and in 108 cases a history of bypass grafting---and
nothing else. They have no comorbidity for any recommendation to match against and, apart from
ejection fraction, no observation to gate on. Roughly 28\% of the evaluated
cohort is therefore clinically far sparser than any real patient with chronic coronary disease
would be.

\begin{figure}[h]
\centering
\begin{tikzpicture}[
    node distance = 4.5mm and 10mm,
    st/.style    = {rectangle, rounded corners=2pt, draw, thick, text width=52mm,
                    align=center, minimum height=7mm, font=\scriptsize},
    onset/.style = {st, fill=black!8},
    side/.style  = {rectangle, rounded corners=2pt, draw, dashed, text width=26mm,
                    align=center, minimum height=6mm, font=\scriptsize},
    ar/.style    = {-latex, thick},
    lb/.style    = {font=\tiny, inner sep=1.5pt}
]
\node[onset]                  (init)  {\textbf{Initial}};
\node[onset, below=of init]   (chance){\textbf{Chance\_of\_IHD} --- \emph{Delay} 1 month};
\node[onset, below=of chance] (start) {\textbf{IHD\_Start} --- \emph{ConditionOnset}\\
                                       SNOMED 414545008 \emph{Ischemic heart disease}};
\node[st, below=of start]     (delay) {\textbf{Delay\_Until\_Symptoms} --- \emph{Delay} 1 month};
\node[st, below=of delay]     (symp)  {\textbf{Symptoms\_Start} $\rightarrow$
                                       \textbf{Dyspnea\_Symptom}\\
                                       chest pain; dyspnoea on exertion};
\node[st, below=of symp]      (enc)   {\textbf{Diagnosis\_Encounter} --- \emph{Encounter}
                                       (wellness)};
\node[st, below=of enc]       (ecg)   {\textbf{ECG / Labs} $\rightarrow$
                                       \textbf{Stress Test}};
\node[st, below=of ecg]       (med)   {\textbf{Medical\_Therapy} --- \emph{Delay} 1 week};
\node[st, below=of med]       (presc) {\textbf{Prescribe\_CHD\_Meds}\\
                                       date-gated cascade: clopidogrel (1997),
                                       simvastatin (1991), metoprolol (1978),
                                       nitroglycerin};
\node[st, below=of presc]     (cabg)  {\textbf{CABG} --- \emph{CallSubmodule}
                                       \texttt{heart/cabg\_sequence}};
\node[st, below=of cabg]      (term)  {\textbf{Terminal}};
\node[side, right=of delay]   (arrest){\textbf{Cardiac\_Arrest} --- \emph{Death}\\
                                       SNOMED 410429000};

\draw[ar] (init)   -- (chance);
\draw[ar] (chance) -- node[lb, right] {$p=\texttt{ihd\_risk}$} (start);
\draw[ar] (start)  -- (delay);
\draw[ar] (delay)  -- node[lb, right] {$0.25$} (symp);
\draw[ar] (symp)   -- (enc);
\draw[ar] (enc)    -- (ecg);
\draw[ar] (ecg)    -- node[lb, right] {$0.99$ at each step} (med);
\draw[ar] (med)    -- (presc);
\draw[ar] (presc)  -- node[lb, right] {$0.9$ via Symptoms\_Continue} (cabg);
\draw[ar] (cabg)   -- (term);
\draw[ar, dashed] (delay)  -- node[lb, above] {$0.0005$} (arrest);
\draw[ar, dashed] (arrest) |- (term);

\draw[ar] (chance.north west) to[out=180, in=180, looseness=2]
          node[lb, left] {else} (chance.south west);
\draw[ar] (delay.north west)  to[out=180, in=180, looseness=2]
          node[lb, left] {$0.7495$} (delay.south west);
\draw[ar] (presc.west)        to[out=180, in=180, looseness=1.5]
          node[lb, left] {$0.1$} (med.west);
\end{tikzpicture}
\caption{The \texttt{stable\_ischemic\_heart\_disease} module as a state machine, at Synthea
commit \texttt{aa0772f}. Only the shaded states determine cohort membership: a patient enters
\texttt{Chance\_of\_IHD} and draws once per simulated month against \texttt{ihd\_risk}, so most
patients loop until death without ever reaching \texttt{IHD\_Start}. Everything below the
shaded region is care pathway, not diagnosis, and runs only for patients who did transition.
Probabilities are those given literally in the module; \texttt{ihd\_risk} is supplied
externally by \texttt{CardiovascularDiseaseModule} and is undefined outside ages 40--79. For
legibility two minor edges are omitted: \texttt{ECG / Labs} and \texttt{Stress Test} each route
a further $0.01$ straight to \texttt{CABG}, bypassing medical therapy.}
\label{fig:sihd-module}
\end{figure}

\subsection{The CCD gate and cohort derivation}
\label{sec:gate}
Cohort membership is therefore established by querying the generated data rather than by
trusting the run configuration. A patient enters the evaluation cohort if the graph contains,
for that patient, a \texttt{:Condition} whose code bridges to the UMLS concept
\texttt{C0151744}---which the terminology layer resolves from SNOMED~CT \texttt{414545008},
\emph{Ischemic heart disease (disorder)}. Concretely, the gate is the traversal
\texttt{(:Patient)$\leftarrow$[:SUBJECT]$-$(:Condition)$-$[:CODE]$\rightarrow$(:Code)
$\leftarrow$[:SAME\_CONCEPT]$-$(:Concept)}, and it is the first clause evaluated in the
applicability query: a patient who fails it is never matched against any recommendation.
The same gate concept is used by the reference standard against which the system is scored,
an automatically generated \emph{silver} standard defined and justified in
Section~\ref{sec:oracle}. Both sides of the concordance evaluation are therefore computed
over an identical cohort.

Table~\ref{tab:cohort-funnel} traces the resulting derivation.

\begin{table}[h]
\centering
\caption{Derivation of the gate-evaluable cohort from the generated corpus.}
\label{tab:cohort-funnel}
\begin{tabular}{lrl}
\hline
\textbf{Stage} & \textbf{Patients} & \textbf{Determined by} \\
\hline
Generated corpus (both runs)      & 5{,}490 & Synthea run configuration \\
\quad disease-enriched run        & 2{,}009 & \texttt{stable\_ischemic\_heart\_disease} \\
\quad general run                 & 3{,}481 & all modules (\texttt{*}) \\
Carrying SNOMED 414545008         &    876  & simulation outcome \\
\textbf{Gate-evaluable cohort}    & \textbf{876} & \textbf{the CCD gate} \\
\hline
\end{tabular}
\end{table}

The 876 figure is thus not a filter applied on top of a larger set of confirmed cases; it is
the number of simulated patients who actually developed the gating diagnosis, or
$16.0\%$ of the corpus.

\subsection{What the gate excludes}
The gate is deliberately narrow, and the consequence should be stated explicitly. The 2023
AHA/ACC guideline defines chronic coronary disease \emph{disjunctively}: prior myocardial
infarction and prior coronary revascularization are CCD-defining alongside a diagnosis of
ischemic heart disease. The gate implemented here accepts only the last of these. A further
67 patients in the corpus carry \emph{History of myocardial infarction} (SNOMED
\texttt{399211009}) or \emph{History of coronary artery bypass grafting} (SNOMED
\texttt{399261000}) without also carrying \texttt{414545008}; on the guideline's own
definition these are CCD patients, which would place the clinically defined cohort at 943.

This gap is not a terminology-bridging failure. All three codes resolve correctly to distinct UMLS concepts: \texttt{414545008} to
\texttt{C0151744}, \texttt{399211009} to \texttt{C1275835}, and \texttt{399261000} to
\texttt{C1275842}. The distinction is semantically correct, because SNOMED~CT models a
\emph{situation with explicit context} (``history of'') as a different entity from the
corresponding \emph{disorder}. The guideline
layer is likewise not at fault: six recommendations already target the history-of-infarction
concept. The limitation is narrower and more specific than either---the gate is expressed as
a \emph{single} concept where the guideline expresses a disjunction, so those six
recommendations can never be reached for the 67 patients, because the gate is evaluated
before recommendation matching begins.

The single-concept gate is retained in this study so that the reported concordance,
agreement, and safety figures all describe one consistently defined cohort. The
gate-evaluable cohort should therefore be read as a conservative subset of the clinical CCD
population---$876$ of $943$, or $92.9\%$---and the residual $7.1\%$ is revisited among the
limitations in Section~\ref{sec:ethics}.
The 30~patients labelled by hand were not drawn at random. Every one of the 876 was
scored for how far it would exercise the briefing. The scoring awards two points per discriminating comorbidity, counted
for at most six of the nine on the allow-list, so contributing at most 12 points. It awards
one point per safety-relevant drug class, counted for at most three of the six, so
contributing at most 3 points. The highest attainable score is therefore 15. The score is then penalised for artefacts of the generator. The heaviest
penalty falls on a left-ventricular ejection fraction outside 20--70\%, which is not
physiological and would invite a reader to criticise Synthea rather than the system. The 30 highest-scoring patients were taken. The
procedure is deterministic and needs no seed: ranks~1--18 reach that ceiling and
ranks~19--30 score 13, with the cut falling inside that tie and settled by identifier order.
The sample is therefore purposive, weighted towards the cohort's most clinically complex
records, and the agreement reported in Section~\ref{sec:goldlabel} should be read as
measured on those cases rather than on an average patient.

\section{Knowledge-Graph Construction Method}
\label{sec:kg-construction}

The knowledge graph is built \emph{EHR-first}: the concept layer is minted from the EHR codes
before the literature corpus is loaded. EHR-to-UMLS resolution depends only on the code
descriptions and the linker's UMLS release, never on which abstracts happen to have been
retrieved, so the set of concepts reachable from patient data---and with it every
\texttt{SAME\_CONCEPT} edge---survives any repetition or widening of the literature search
unchanged. This matters because the CCD gate of Section~\ref{sec:gate} traverses precisely
those edges. Had the concept layer been minted from the literature instead, a different corpus
would give a different concept set and therefore a different cohort, making the denominator of
every reported metric a function of the literature search. Literature-derived concepts are
merged onto the same \texttt{umlsCui} node key afterwards, enriching the layer without
redefining it. The construction proceeds as follows.

\begin{enumerate}
    \item \textbf{FHIR-to-graph loading.} Each FHIR resource becomes a node, each reference
    becomes a typed relationship, and every terminology code is hoisted into a shared
    \texttt{:Code} node keyed by (system, code) so that all occurrences of a code refer to
    one node. Multi-valued observations (e.g.\ blood pressure) are represented as
    component nodes to avoid data loss. Billing/administrative resources are skipped, and
    only an allow-list of clinically relevant LOINC observations is loaded.
    \item \textbf{Terminology bridge (EHR-to-UMLS).} The deduplicated EHR codes are exported
    from the graph and mapped to UMLS \cite{umls} concept unique identifiers (CUIs) using a
    three-pass linker (exact alias, candidate generation, and NLP-based disambiguation
    \cite{scispacy}), producing the
    concept layer connected to the shared \texttt{:Code} nodes via \texttt{SAME\_CONCEPT}
    edges.
    \item \textbf{Guideline layer.} The 2023 AHA/ACC CCD guideline is hand-curated into
    81 recommendations, each linked to the concepts it recommends,
    contraindicates, or targets. The edge type is derived from the Class of Recommendation
    (COR): COR~3 (Harm / No-Benefit) yields a \emph{contraindicates} edge, and COR~1/2a/2b
    yields a \emph{recommends} edge. Disjunctive-target recommendations are tagged so the
    query layer joins their targets with OR rather than AND.
    \item \textbf{Literature layer.} A guideline-seeded PubMed corpus (the guideline's cited
    trials plus their similar articles) is processed with scispaCy clinical NER to extract
    entity mentions, which are linked to the same UMLS concepts. Drug--condition
    associations are materialised between \texttt{:Code} nodes as co-occurrence edges ranked
    by pointwise mutual information (PMI) \cite{pmi}, with a minimum article-count floor.
    The corpus comprises 4{,}801 distinct abstracts: the 274 trials the guideline
    itself cites, together with 4{,}793 similar articles retrieved from them. The two sets
    overlap heavily---266 of the seeds resurface among their own similar-article results, and
    each abstract is stored once---so the union is smaller than the sum. A drug--condition pair
    must be co-mentioned in at least three distinct abstracts before it is eligible for
    ranking.
\end{enumerate}

\subsection{Graph data model}
The loader maps FHIR to the property graph as follows. Each resource becomes a node labelled
\texttt{:Resource:<ResourceType>}, with its scalar fields flattened onto the node. A field
containing a reference becomes a typed relationship whose name is derived from the FHIR
element path; for example, \texttt{Observation.subject} yields a \texttt{SUBJECT} edge. The
loader resolves the three reference forms Synthea emits: \texttt{urn:uuid:} identifiers,
conditional references (\texttt{Type?identifier=system\textbar value}), and contained
(\texttt{\#fragment}) references.

The shared \texttt{:Resource} label is plumbing rather than meaning. It says only
that a node can be found by its \texttt{uid}; it claims nothing about whether that node is a
Resource in the FHIR sense. The label exists because of what the loader knows, and when.
Every reference form collapses to a bare identifier: \texttt{Type/id} drops its prefix, and
\texttt{urn:uuid:} never had one. By the time an edge is written, the loader therefore has
the target's identifier but not its type. It therefore needs one label that every possible target
shares, and this is that label. The same label holds the single \texttt{uid} uniqueness
constraint, which the database backs with an index, so following a reference costs one index
lookup instead of a scan.

Two consequences follow, both intentional. Component nodes carry the \texttt{:Resource}
label even though \texttt{Observation.component} is a BackboneElement in FHIR rather than a
Resource, because edges must be able to resolve to them by identifier. Conversely,
\texttt{:Code} nodes are the deliberate exception: they are keyed by
(\texttt{system}, \texttt{code}) rather than by \texttt{uid}, are never the target of a
resource reference, and therefore carry no \texttt{:Resource} label. Every terminology coding is hoisted into a shared node
\texttt{:Code\{system, code, display\}}, keyed by (\texttt{system}, \texttt{code}) under a
uniqueness constraint, so that every occurrence of a code---across conditions, procedures,
medications, and observations---refers to a single node; this shared \texttt{:Code} node is
the anchor to which the UMLS concept layer attaches. A Coding is distinguished from a
Quantity (both carry a \texttt{system} and a \texttt{code}) by the presence of a
\texttt{value}, which prevents units of measurement from being promoted to clinical codes.
Multi-valued observations---notably the blood-pressure panel---are expanded into
\texttt{:Component} nodes so that the systolic and diastolic values are both preserved. All
writes use \texttt{MERGE}, which makes the load idempotent and re-runnable; billing and
provenance resources are skipped, and only the allow-listed LOINC observations are loaded.

\subsubsection*{Why several terminologies coexist}
\label{sec:multi-vocab}
The loaded graph contains 26 distinct code systems. This is a property of FHIR
rather than a quirk of the generator: the specification binds each coded element to a value
set, and those bindings name particular vocabularies. Synthea declares that conformance
explicitly---smoking-status observations carry a \texttt{meta.profile} reference to
\texttt{us-core-smokingstatus}---so the vocabularies are those the US~Core profiles make
mandatory, not free choices of the generator. Each vocabulary answers one question and
does not attempt the others. SNOMED~CT \cite{snomedct} carries clinical meaning---conditions,
procedures, encounter and device types---as a description-logic ontology with computable
subsumption. RxNorm \cite{rxnorm} carries medications, normalising drug naming across
proprietary vocabularies at explicit levels of abstraction. LOINC \cite{loinc} carries
observations, identifying \emph{what was measured} under a formal multi-axis model. CVX
carries vaccines. A further group of HL7 code systems carries FHIR's own administrative
values---clinical status, category, encounter class---which are structural metadata rather
than clinical terminology, and which account for roughly half of the distinct systems present.

The division is directly visible in the loaded graph. Observations reach LOINC through
\texttt{CODE} edges (242{,}217) but reach SNOMED~CT through \texttt{VALUE} edges
(47{,}259): LOINC names the question, and SNOMED~CT names the answer wherever the
answer is itself a coded concept rather than a number. Tobacco smoking status is the clearest
instance---LOINC \texttt{72166-2} identifies the observation, and its value is a SNOMED~CT
concept.

Two design consequences follow, and both are load-bearing. First, \texttt{:Code} nodes must be
keyed by (\texttt{system}, \texttt{code}) rather than by code alone, because bare code strings
collide across vocabularies. Second, and more consequentially, the vocabularies are disjoint
namespaces maintained by different organisations on independent release cycles, and no edge
runs between them in any source terminology. A single guideline recommendation routinely spans
several at once: statin therapy for a patient with chronic coronary disease names an RxNorm
ingredient and a SNOMED~CT disorder, and once a lipid threshold is attached, a LOINC
measurement as well. Nothing in the patient data connects the three. The UMLS concept layer is
therefore not a convenience but a structural necessity---it is the only artefact that assigns
them a common identifier, and without it the guideline layer cannot reach the patient layer at
all.

A third consequence bears specifically on this work's contribution. LOINC separates the
measurement from its result, leaving the numeric value in \texttt{valueQuantity}. Had ejection
fraction been recorded the way SNOMED~CT would record it---as a finding such as ``left
ventricular systolic dysfunction''---there would be no number to compare against a threshold,
and the observation gates of Section~\ref{sec:dss-method} could not be written at all. The
observation-gate gain reported in Chapter~IV rests on this property of LOINC specifically, not
on structured encoding in general. Synthea's tabular export has that property too, since it
also carries LOINC codes with their numeric values (Section~\ref{sec:csv-gating}).

\subsubsection*{What a flattened tabular export loses}
\label{sec:flattening}
Chapter~I motivates the FHIR-native design by contrast with flattened tabular exports. Because
Synthea emits the same cohort in both formats, that contrast can be shown rather than asserted.
Table~\ref{tab:flattening} states it for the blood-pressure panel.

\begin{table}[h]
\centering
\caption{The blood-pressure panel in the CSV export and in the FHIR graph, same cohort.}
\label{tab:flattening}
\begin{tabular}{lrl}
\hline
\textbf{LOINC} & \textbf{CSV rows} & \textbf{Graph} \\
\hline
85354-9 blood-pressure panel & 0        & one \texttt{:Observation} per reading \\
8480-6 systolic              & 52{,}583 & \texttt{:Component} of that Observation \\
8462-4 diastolic             & 52{,}583 & \texttt{:Component} of that Observation \\
\hline
\end{tabular}
\end{table}

The panel does not survive the export: it has no row of its own, and its two members become
independent rows, so recovering ``this patient's pressure was 121/80'' means reassembling them.
The graph holds 105{,}166 \texttt{:Component} nodes, which is $2 \times 52{,}583$---every member
preserved, each still attached to the panel that grouped it.

How expensive that reassembly is, is an empirical question, and Section~\ref{sec:csv-gating}
answers it by rebuilding the observation gate from the export. The answer belongs here too, so that this contrast is not read as stronger than it is. The
export carries an \texttt{ENCOUNTER} column populated on every observation row. The two
members of one reading can therefore be brought back together through an explicit foreign
key rather than through timestamps assumed to agree. In the
evaluation cohort, 17{,}177 pairs reassemble completely, leaving no row unpaired. What is lost is
therefore the panel's identity as a unit, not the recoverability of its values.

A second difference is subtler, and its consequence is latent rather than realised. Smoking
status appears in the CSV as 47{,}259 rows carrying only a display
string---35{,}541 \texttt{Never smoked tobacco (finding)}, 11{,}587
\texttt{Ex-smoker (finding)}, and 131 \texttt{Smokes tobacco daily
(finding)}---with \texttt{TYPE} recorded as \texttt{text}. The FHIR resource carries the same
string inside a \texttt{valueCodeableConcept}, but alongside \texttt{system}
\texttt{http://snomed.info/sct} and \texttt{code} \texttt{266919005}. The export keeps the
label and drops the code; the graph keeps both.

Three qualifications are owed, because on this corpus the difference is not merely smaller than
it appears but entirely inert. First, the concept bridge exports codes from conditions, medications, care-plan categories
and procedures, but not from observation values, so \texttt{266919005} never becomes a
\texttt{:Concept} node. Second, the smoking gate reads the value node's \texttt{display}
property rather than its \texttt{code}, so smoking status is matched by string on both
sides. Third, no patient in the evaluation cohort is a current smoker at all. The five
smoking-gated recommendations therefore record zero true positives, false positives and
false negatives, as Section~\ref{sec:ablation-class} reports.

Smoking is therefore a poor example for the claim that codes survive the
export, and it is not used as one here. It is kept because the \emph{kind} of loss it shows is
real and general: a flat table has nowhere to record which coding system a particular value
came from. On this corpus that loss happens to cost nothing. What actually carries the argument
for FHIR is the blood-pressure panel above, whose component values are read as numbers and do
change the results.

The vocabularies themselves are not lost, and it is worth being precise about why. Each
\texttt{CODE} column draws on a vocabulary fixed by the file rather than declared on the row:
\texttt{conditions.csv} and \texttt{procedures.csv} are SNOMED~CT, \texttt{medications.csv} is
RxNorm, \texttt{observations.csv} is LOINC, \texttt{immunizations.csv} is CVX, and the data
dictionary states each of these explicitly \cite{synthea_csv_dict}. A \texttt{SYSTEM} column
appears in precisely the two files whose documentation hedges the vocabulary as
\emph{typically} SNOMED~CT, and is absent from the three whose vocabulary the dictionary states
outright. The asymmetry is therefore principled, and at the level of whole columns the export
loses nothing.

What a fixed-per-file vocabulary cannot express is a code whose system varies \emph{within} a
column---which is exactly the observation \texttt{VALUE} case above. That single column holds a
number for LDL cholesterol and a SNOMED~CT concept for smoking status, so there is no
consistent system for a companion column to declare, and the code is dropped in favour of the
display string. FHIR attaches \texttt{system} to each individual coding and so faces no such
constraint. The loss is a consequence of the tabular shape, not of careless schema design.

One further loss concerns precision rather than structure. The export writes observation values
to one decimal place, where the FHIR resource keeps them whole. The difference is bounded at
$0.05$ and can only matter where a value happens to sit right at the edge of a guideline
threshold. Section~\ref{sec:csv-gating} counts how often that happens on this cohort.

The genuine structural omission is \texttt{status}. No table carries it: the FHIR resource records
\texttt{"status": "final"} and the CSV has nowhere to put it. That same absence in
\texttt{medications.csv}, which offers \texttt{START}, \texttt{STOP} and \texttt{DISPENSES} but
no status, is precisely what Section~\ref{sec:chronicity} must work around.
Both totals below come entirely from the general run: the disease-enriched run records no
blood-pressure and no smoking-status observation at all.

\subsubsection*{Why observations are allow-listed}
Observations are the dominant scale term in a FHIR corpus: they account for roughly
38\% of all resources, and a single patient in this cohort can carry several
thousand of them. Loading every one would multiply the graph several-fold with routine
measurements that no guideline recommendation consults---repeated urinalysis strips, vital
signs from unrelated encounters, panel members never referenced by the encoded logic. The
loader therefore applies an allow-list of nine LOINC codes covering exactly the
values the guideline logic turns on: ejection fraction, body-mass index, LDL and total
cholesterol, triglycerides, glycated haemoglobin, smoking status, and the blood-pressure
panel. Systolic and diastolic pressures need no separate entry because they are carried as
components inside the panel and are captured with it.

This is a decision about \emph{scale}, and it is worth separating from the observation-gating
mechanism described in Section~\ref{sec:dss-method}, which is a decision about
\emph{expressiveness}. The allow-list controls which observations exist in the graph at all;
observation gating controls how a recommendation that depends on a numeric value is decided
once they are there. Widening the allow-list would enlarge the graph without changing how
gating works, and narrowing it would starve the gates without altering their logic.
The allow-list is derived from the encoded guideline logic rather than imposed on it: a
code is admitted precisely because some recommendation reads it. No recommendation is
therefore excluded from the evaluation for want of an observation---the list bounds how much
of the corpus is loaded, not which recommendations can be assessed.

\subsection{Two views of the schema}
The loaded graph is described here from two complementary angles, because the complete schema
and the part of it that the recommendation layer actually traverses differ greatly in size.
The \emph{full} schema is the set of all label--relationship--label triples present in the
database: 24 node labels---excluding the structural \texttt{:Resource} label, which by
construction attaches to everything except \texttt{:Code}---joined by 95 distinct triples. It
is reproduced in Appendix~\ref{app:schema}. The \emph{usable} schema is the subgraph the
queries traverse: 11 labels and 11 relationship types. Twelve clinical labels are loaded
faithfully but read by no query reported in this thesis---\texttt{:AllergyIntolerance},
\texttt{:CarePlan}, \texttt{:CareTeam}, \texttt{:Device}, \texttt{:Encounter},
\texttt{:ImagingStudy}, \texttt{:Immunization}, \texttt{:Location},
\texttt{:MedicationAdministration}, \texttt{:Organization}, \texttt{:Practitioner} and
\texttt{:PractitionerRole}---and they remain available for future work. The thirteenth
unqueried label, \texttt{:Guideline}, is left out for a different reason: it carries the
guideline's provenance through
\texttt{(:Guideline)$-$[:CONTAINS]$\rightarrow$(:Recommendation)}, and the recommendation
queries enter the guideline layer at \texttt{:Recommendation} itself, so the parent node is
never needed at query time.

The usable subgraph was not chosen by judgement. It is the set of labels and relationship
types appearing literally in the recommendation, observation-gate, briefing, and chronicity
modules, and every edge in Figure~\ref{fig:graph-model} was verified to exist in the loaded
database rather than inferred from the loader source.

The shared \texttt{:Code} node is
the pivot of the whole design: every clinical resource type converges on it, and it is the
single attachment point for the UMLS concept layer, which in turn carries the guideline and
literature layers. This is what allows a patient to be connected to a recommendation and to
its supporting evidence by traversal alone, without any string matching at query time. Note
the direction of the bridge: the edge runs
\texttt{(:Concept)-[:SAME\_CONCEPT]->(:Code)}, from concept to code, so a traversal starting
at the patient reaches concepts by following it in reverse.

\begin{figure}[htbp]
\centering
% Scaled to the measure: the natural width is a few mm over \textwidth.
\resizebox{\textwidth}{!}{%
\begin{tikzpicture}[
    node distance = 3.5mm and 10mm,
    res/.style  = {rectangle, rounded corners=2pt, draw, text width=27mm,
                   align=center, minimum height=7mm, font=\scriptsize},
    hub/.style  = {rectangle, rounded corners=2pt, draw, very thick, text width=21mm,
                   align=center, minimum height=9mm, font=\scriptsize, fill=black!12},
    lay/.style  = {rectangle, rounded corners=2pt, draw, text width=23mm,
                   align=center, minimum height=9mm, font=\scriptsize},
    pat/.style  = {rectangle, rounded corners=2pt, draw, very thick, text width=17mm,
                   align=center, minimum height=9mm, font=\scriptsize, fill=black!6},
    lbl/.style  = {font=\tiny, inner sep=1pt},
    ar/.style   = {-latex}
]
% EHR resources
\node[res] (c)  {\texttt{:Condition}};
\node[res, below=of c] (p)  {\texttt{:Procedure}};
\node[res, below=of p] (m)  {\texttt{:MedicationRequest}};
\node[res, below=of m] (md) {\texttt{:Medication}};
\node[res, below=of md] (o) {\texttt{:Observation}};
\node[res, below=of o] (cm) {\texttt{:Component}};

% patient anchor
\node[pat, left=10mm of p, yshift=-8mm] (pt) {\texttt{:Patient}};

% terminology bridge and upper layers
\node[hub, right=of p, yshift=-10mm] (code) {\texttt{:Code}\\(system, code)};
\node[lay, right=of code] (con) {\texttt{:Concept}\\(UMLS CUI)};
\node[lay, above right=6mm and 10mm of con] (rec) {\texttt{:Recommendation}\\(COR)};
\node[lay, below right=6mm and 10mm of con] (abs) {\texttt{:Abstract}\\(PMID)};

% SUBJECT: clinical resources anchor on the patient
\draw[ar] (c) -- (pt);
\draw[ar] (p) -- (pt);
\draw[ar] (m) -- (pt);
\draw[ar] (o) -- (pt);
\node[lbl, above=0.5mm of pt] {\texttt{SUBJECT}};

% CODE / MEDICATION / COMPONENT
\draw[ar] (c)  -- (code);
\draw[ar] (p)  -- (code);
\draw[ar] (o)  -- (code);
\draw[ar] (cm) -- (code);
\draw[ar] (md) -- (code);
\draw[ar] (m)  -- (code);
\draw[ar] (o)  -- (cm);
\draw[ar] (m)  -- (md);

% bridge runs Concept -> Code
\draw[ar] (con) -- node[lbl, above, sloped] {\texttt{SAME\_CONCEPT}} (code);
\draw[ar] (rec) -- (con);
\draw[ar] (abs) -- (con);
\draw[ar] (rec) -- node[lbl, right] {\texttt{CITES}} (abs);

\node[lbl, above=1mm of c]   {EHR layer};
\node[lbl, above=1mm of rec] {guideline layer};
\node[lbl, below=1mm of abs] {literature layer};
\end{tikzpicture}}
\caption{The \emph{usable} schema: the 11 labels and 11 relationship types traversed by the
recommendation queries. Clinical resources anchor on \texttt{:Patient} via \texttt{SUBJECT}
and converge on the shared \texttt{:Code} node; \texttt{:Concept} bridges \emph{to}
\texttt{:Code} via \texttt{SAME\_CONCEPT}, carrying the guideline and literature layers.
Medications resolve by two paths---inline to \texttt{:Code}, or by reference through
\texttt{:Medication}---both of which the queries follow. A
\texttt{CO\_OCCURS\_IN\_LITERATURE} edge, omitted here for legibility, additionally joins
\texttt{:Code} to \texttt{:Code}. The complete 95-triple schema is given in
Appendix~\ref{app:schema}.}
\label{fig:graph-model}
\end{figure}

\section{Recommendation Method}
\label{sec:dss-method}

Given a patient, the recommendation layer performs four deterministic operations against
the graph:

\begin{enumerate}
    \item \textbf{Applicable recommendations.} A recommendation is applicable if the patient
    belongs to the CCD cohort (the baseline gate) \emph{and} the patient matches all of the
    recommendation's condition targets. Condition targets are the further diagnoses, beyond
    the cohort gate itself, that a recommendation may require, and they are counted only
    where they bridge to a code the EHR actually uses. The traversal is gate-plus-AND, with
    OR for disjunctively tagged recommendations.
    Section~\ref{sec:disc-class} formalises the term.
    \item \textbf{Observation gating.} Recommendations whose preconditions depend on numeric
    observations (e.g.\ ejection fraction, LDL, BMI, triglycerides, blood pressure, smoking
    status) are decided by evaluating the guideline predicate on the patient's measured
    values read from the graph, using the same thresholds as the reference oracle
    (Section~\ref{sec:oracle}). This is the mechanism through which the FHIR-native
    observations improve precision.
    \item \textbf{Contraindication safety.} The patient's medications are checked against
    the recommendations that contraindicate them, in three variants. The unary variant flags
    if the patient is on the drug. The population-gated variant flags only if the patient
    also meets the contraindication's population condition. The chronicity-gated variant,
    used for the chronic-NSAID recommendation, flags only if the drug qualifies as chronic
    use. The chronicity gate used inside the recommendation layer is computed from the graph alone.
    It reads \texttt{MedicationRequest.status}, which FHIR records as a category
    (\texttt{active} or \texttt{completed}) rather than as a date, together with the count
    of separate prescribing episodes. It therefore exercises only the subset of the
    chronicity definition that survives the FHIR export. Because no stop date is exported, the two arms
    of that definition expressed as durations cannot be evaluated from the graph at all;
    Section~\ref{sec:chronicity} sets out the asymmetry in full.
    \item \textbf{Literature support and briefing.} Drug options linked to the patient's
    conditions in the literature layer are ranked by PMI, and all of the above facts are
    assembled into a structured payload that a large language model phrases into a
    clinician-facing briefing. The model never queries the graph and never makes a clinical
    decision; it only re-words the retrieved, provenance-carrying facts.
\end{enumerate}

\subsection{Query design and the observation-gate procedure}
Applicability is evaluated by a single traversal. Starting from the patient, the query first
requires membership in the CCD cohort (the baseline gate concept). For each candidate
recommendation, it then requires that the patient match \emph{all} of the recommendation's
non-anchor condition targets that bridge to the EHR. This is implemented as a logical AND:
the number of matched targets must equal the number of bridgeable targets. Recommendations
tagged as disjunctive take an OR branch instead, so one matched target suffices. Because the
medication choice type in FHIR may be expressed either inline or by reference, both
medication paths (the inline \texttt{:Code} and the \texttt{:Medication}$\to$\texttt{:Code}
reference) are resolved wherever medications are matched, so that referenced drugs are not
silently dropped.

\subsubsection*{How a recommendation discriminates}
\label{sec:disc-class}
Not every recommendation narrows the cohort. Once a patient has passed the CCD gate, some recommendations apply to that patient
automatically, with nothing further to test. Others apply only to patients who meet a
further condition. Which of the two a recommendation is, and by what mechanism it selects,
governs how the results of Chapter~IV should be read. Two mechanisms are available for
narrowing, and they are independent. A recommendation may carry a \emph{condition target}
beyond the gate concept that bridges to the EHR, in which case the traversal can exclude
patients who do not match it. Separately, it may carry a \emph{numeric precondition} over an
observation, in which case the observation gate can exclude patients whose measured value
fails the predicate. Crossing the two yields four classes, summarised in
Table~\ref{tab:disc-class}.

\begin{table}[h]
\centering
\caption{Classification of the non-contraindication recommendations by the mechanism
available for narrowing the cohort.}
\label{tab:disc-class}
\begin{tabular}{lccl}
\hline
\textbf{Class} & \textbf{Target?} & \textbf{Gated?} & \textbf{Recommendations} \\
\hline
Universal        & no  & no  & 38 \quad (NUTRI, PA, WM:1, APT:1--2 \ldots) \\
Value-gated      & no  & yes & \phantom{0}5 \quad (BB:1, BB:2, BPM:1, RAAI:1, RAAI:2) \\
Concept-matched  & yes & no  & 15 \quad (APT:11--13, APT:4--5 \ldots) \\
Both             & yes & yes & 10 \quad (LIPID:9, SGLT2:4, SGLT2:6, TP:2 \ldots) \\
\hline
\textbf{Total}   &     &     & \textbf{68} \\
\hline
\end{tabular}
\end{table}

The \emph{universal} class is the largest. The nutrition, physical-activity and
weight-management recommendations belong to it because the guideline directs them at every
patient with CCD; a system that fires them for every cohort member is behaving correctly.
\emph{Concept-matched} recommendations discriminate by traversal---those addressing patients
who also require oral anticoagulation apply only where that indication is coded.
\emph{Value-gated} recommendations discriminate only by predicate. \texttt{BB:1} and
\texttt{BB:2}, for instance, turn on ejection fraction, which is a measured number and not a
diagnosis. No traversal can therefore separate the patients they apply to. The \emph{both} class is subject to
each mechanism in turn.

Two consequences follow, and the first is easy to state incorrectly. It is tempting to say
that a universal recommendation cannot lower precision, on the grounds that it applies to
every cohort member by definition. That is false, and Section~\ref{sec:ablation-class}
measures how false: the universal class generates the large majority of the system's false
positives. The reason is that ``universal'' here is a statement about the \emph{encoding},
not about the guideline. It means the recommendation carries no condition target that both
was encoded and bridges to an EHR code---not that the guideline text imposes no
precondition. Where the guideline does impose one that was never represented in the
knowledge layer, the system has no means of withholding the recommendation, and the
reference standard, which reads the guideline logic from the raw record, marks the
difference as a false positive. The dominant limit on precision is therefore the
\emph{coverage of the encoded guideline}, not the quality of the matching applied to it.

Second, and this one is sharper. The observation-gate ablation of
Section~\ref{sec:eval-method} can only touch the 15 recommendations that are gated at all:
the five value-gated ones and the ten in the \emph{both} class. The other 53 are scored
exactly the same way whether gating is switched on or off. So whatever difference the ablation
shows has to come from those 15, and from nowhere else. That is a prediction rather than an
assumption, which is why Chapter~IV reports concordance class by class---so that it can be
checked. It holds exactly: for the ungated classes, the two arms agree down to the individual
counts.

\subsubsection*{Why observation gating is a separate mechanism}
Observation-dependent recommendations cannot be resolved by the traversal just described, for
a reason of expressiveness rather than of scale. The gate-plus-AND query matches
\emph{concepts}: it asks whether a patient's record contains a code that bridges to the
concept a recommendation targets. But a precondition such as ``LDL-C below 70~mg/dL'' or
``ejection fraction at or below 40\%'' is not a concept---it is a predicate over a
\emph{value}. No pattern over nodes and edges evaluates it, because what must be compared is
a number carried as a property, against a threshold stated in the guideline text and not
present in the graph at all.

Observation-dependent recommendations are therefore decided separately: the observation-gate
module reads the patient's aggregated observation features from the graph (worst-ever
laboratory values, most-recent blood pressure and smoking status) and evaluates the guideline
predicate for each such recommendation using the same thresholds as the reference standard.
The final applicable set is assembled from two halves. Recommendations without an
observation gate keep whatever the graph's gate-plus-AND traversal decided. Gated ones are
decided by their predicate instead, and that decision \emph{replaces} the graph's rather than
being intersected with it. The difference shows up wherever the guideline states a precondition
as a list of alternatives. One recommendation applies if the patient has hypertension, diabetes, chronic kidney
disease, \emph{or} a reduced ejection fraction. Matching all four at once would drop every
patient who has only one of them. Letting the predicate decide brings those patients back.

The grounded briefing is produced in four stages. Stage~(i) is deterministic retrieval of
the patient snapshot, applicable recommendations, safety flags, and literature support via
the queries above. Stage~(ii) assembles these into a structured facts payload in which every
item carries its provenance: a recommendation identifier, a class of recommendation, and
PMIDs. Stage~(iii) is a single language-model call that phrases the facts; it must justify
each recommendation only from its retrieved trigger condition, and it may not assert
preconditions the data do not support. Stage~(iv) outputs the facts and the briefing as
separate artifacts so that their citation-level faithfulness can be checked mechanically. Listing~\ref{lst:applicable-core} gives the core of that query; the full text is
reproduced in Appendix~\ref{app:queries}.
\begin{lstlisting}[caption={Core of the applicable-recommendation query: the CCD gate, then
target matching under the recommendation's own conjunction/disjunction flag.},
label={lst:applicable-core}]
MATCH (p:Patient {uid:$pid})
WHERE EXISTS { (p)<-[:SUBJECT]-(:Condition)-[:CODE]->(:Code)
                 <-[:SAME_CONCEPT]-(:Concept {umlsCui:$gate}) }
MATCH (rec:Recommendation) WHERE NOT rec.cor STARTS WITH '3'
/* ... collect each recommendation's non-anchor targets as extraConds,
   and those the patient actually meets as matched ... */
WHERE size(extraConds) = 0
   OR (coalesce(rec.targetLogic,'AND') = 'OR'  AND size(matched) >= 1)
   OR (coalesce(rec.targetLogic,'AND') = 'AND' AND size(matched) = size(extraConds))
RETURN rec.recommendationId AS recId
\end{lstlisting}

Figure~\ref{fig:briefing-pipeline} shows this arrangement. The essential property is that
the language model sits \emph{downstream} of retrieval and has no path back to the graph.
It has no way to obtain a recommendation outside what retrieval supplied, so any
recommendation outside that set must be fabricated. That is precisely what makes the check of
citable items decidable. The arrangement does not guarantee that every sentence the model
writes is semantically faithful (Section~\ref{sec:faithfulness}).

\begin{figure}[htbp]
\centering
\begin{tikzpicture}[
   node distance = 8mm,
   box/.style = {rectangle, rounded corners=2pt, draw, text width=24mm, align=center,
                 minimum height=13mm, font=\scriptsize},
   ar/.style  = {-latex, thick}
]
\node[box] (g) {Knowledge graph\\[1pt]\emph{deterministic Cypher retrieval}};
\node[box, right=of g] (f) {Facts payload\\[1pt]\emph{recId, COR, PMIDs}};
\node[box, right=of f] (l) {Language model\\[1pt]\emph{phrasing only}};
\node[box, right=of l] (b) {Briefing\\[1pt]\emph{main + details}};
\node[box, below=10mm of l, fill=black!8] (chk) {Faithfulness check};

\draw[ar] (g) -- (f);
\draw[ar] (f) -- (l);
\draw[ar] (l) -- (b);
\draw[ar] (f) |- (chk);
\draw[ar] (b) |- (chk);
\end{tikzpicture}
\caption{Grounded-briefing pipeline. Retrieval is deterministic; the language model only
phrases the retrieved facts, and the faithfulness check compares its output back against
the same facts payload.}
\label{fig:briefing-pipeline}
\end{figure}

\subsection{Medication chronicity: FHIR status versus CSV durations}
\label{sec:chronicity}
One recommendation---avoid chronic NSAID use in chronic coronary disease---depends not on
\emph{whether} a drug was prescribed but on \emph{how long} it was taken, and it exposes an
asymmetry between the two faces of the same synthetic record. The FHIR \texttt{MedicationRequest} carries an \texttt{authoredOn} date and a categorical
\texttt{status} (\texttt{active} versus \texttt{completed}) but \emph{no} stop date. The
paired Synthea CSV export, conversely, carries explicit \texttt{START} and \texttt{STOP}
dates but \emph{no} status field. The two encode the same underlying fact through different
projections: an empty CSV \texttt{STOP} (the course is still running) corresponds one-to-one
to a FHIR \texttt{status} of \texttt{active}. Neither projection is a superset of the other.

\subsubsection*{Operationalising ``chronic''}
The guideline does not define the term it turns on. Recommendation
\texttt{CCD-2023:APT:10} (COR~3-Harm, LOE~B-R) reads in full: \emph{``In patients with CCD,
chronic nonsteroidal anti-inflammatory drugs should not be used because of increased
cardiovascular and bleeding complications.''} It specifies no duration, no number of
courses, and no dosing pattern. Any implementation must therefore supply an operational
definition of its own, and that definition is a research decision to be declared rather
than a fact to be cited. This is itself a finding: a recommendation that is unambiguous to
a human reader is not directly executable, and the gap must be closed by the implementer.

No single published definition covers the case, but the published thresholds are not
arbitrary either: they cluster. At the short end, a systematic review of short-term NSAID
safety adopts ``ten days or fewer'' and notes that the United States Food and Drug
Administration advises clinician consultation when use extends beyond ten days
\cite{aminoshariae2016}, and an osteoarthritis expert review describes usual initial courses
as seven to ten days with reassessment by about three weeks
\cite{magni2021}. Beyond that, explicit thresholds appear at three points. A
pharmacovigilance study treats more than 21 consecutive days of NSAID
exposure as potentially inappropriate or overexposed use \cite{ussai2015}. A
population study counts NSAID use from self-reported consumption for at least
30 days \cite{hessami2025}. A cross-sectional study of
prescribing in the elderly states directly that a patient may be considered a chronic NSAID
user on consuming NSAIDs for at least three months
\cite{abdu2020}. Others avoid a day count altogether, classifying chronic use
by medication possession ratio \cite{campbell2022} or by percentage of days
supplied \cite{kikuchi2021}.

The definition adopted in this work has four arms. Two of its thresholds sit inside the
published cluster just described and are chosen for that reason; the remaining one is a
declared parameter with no direct precedent:

\begin{enumerate}
    \item[(i)] an ongoing (active) course;
    \item[(ii)] cumulative exposure exceeding 90~days;
    \item[(iii)] three or more separate courses;
    \item[(iv)] a single course exceeding 30~days.
\end{enumerate}

Arm~(ii), cumulative exposure beyond 90~days, corresponds to the three-month cutoff stated
explicitly as chronic use by Abdu \textit{et al.} \cite{abdu2020}, the
clearest single threshold in this literature. Arm~(iv), a single course beyond 30~days, sits
between the two nearest published boundaries. The first is the 21 consecutive days
that Ussai \textit{et al.} treat as already inappropriate
\cite{ussai2015}. That is the only threshold in the group expressed in
\emph{consecutive} days, and therefore the closest in construction to a single continuous
course. The second is the 30-day mark used by Hessami \textit{et al.}
\cite{hessami2025}. Arm~(iii), three or more separate courses, has no direct
precedent; the coverage-ratio methods \cite{campbell2022}\cite{kikuchi2021}
capture the same intuition of repeated rather than one-off use, but not in a form that
transfers to a count of prescription episodes. It is declared here as a parameter of this
study.

Two qualifications belong with this. The first is that what agrees with the
published work is the \emph{size} of the numbers, not the way they were arrived at. Hessami
\textit{et al.} measure what patients report having taken over time, and Abdu \textit{et al.}
classify patients rather than courses, while the arms here are computed from the
\texttt{START} and \texttt{STOP} dates on prescriptions. The published thresholds justify the
numbers chosen; they do not show that the same thing is being counted.

The second is that the rule is deliberately cautious, and it is better to say so than to call
it accurate. It is \emph{specific rather than sensitive}: it would sooner miss a chronic user
than flag someone who is not one. A patient at fifteen days is not caught, even though fifteen
days is already past the point where unsupervised short-term use is discouraged
\cite{aminoshariae2016} and is approaching the threshold Ussai \textit{et al.} would flag. For
a COR~3 (Harm) recommendation, missing cases is the failure accepted here, so that every flag
the system does raise is beyond argument. The graph-side gate described below makes the same
trade, so this recommendation is treated the same way throughout.

Critically, the \emph{finding} reported in this thesis does not depend on these numbers.
Whatever thresholds are chosen, arms~(ii) and~(iv) are expressed in durations and arms~(i)
and~(iii) are not. The claim that the FHIR projection supports arms~(i) and~(iii) but not
arms~(ii) and~(iv) is a statement about what the export preserves. It says nothing about
where the cutoffs fall. Only the
numeric recall reported in Chapter~IV would move if the thresholds were changed.

Of the four arms, only~(i) and~(iii) are computable from
the graph---arm~(i) from the categorical \texttt{status}, and arm~(iii) from the count of
NSAID \texttt{:MedicationRequest} nodes. The duration-based arms~(ii) and~(iv) require stop
dates and are therefore invisible to a purely FHIR-native reader. The reference standard
(Section~\ref{sec:oracle}) evaluates all four arms from the CSV export, which is still raw
synthetic data and never the CRS output. The graph-side gate embedded in the recommendation
and briefing layers, by contrast, evaluates only arms~(i) and~(iii). The gap between the two is not an
implementation shortcut but a direct consequence of what the FHIR export preserves; its cost
in safety-flag recall is quantified in Chapter~IV. Figure~\ref{fig:chronicity} sets out the
four arms and marks the two that survive the FHIR projection. The rule is implemented in
\texttt{silver\_oracle.py} over eight NSAID ingredients---ibuprofen, naproxen, celecoxib,
diclofenac, meloxicam, indomethacin, ketorolac and piroxicam---matched on the medication
description.

\begin{figure}[htbp]
\centering
\begin{tikzpicture}[
  node distance = 3mm,
  arm/.style = {rectangle, rounded corners=2pt, draw, text width=95mm, align=left,
                minimum height=8mm, font=\scriptsize, inner sep=4pt},
  hit/.style = {arm, fill=black!12, very thick},
  miss/.style= {arm, dashed}
]
\node[hit]  (a1) {(i)~~Ongoing course \hfill computable: \texttt{status = active}};
\node[miss, below=of a1] (a2)
    {(ii)~~Cumulative exposure $>$ 90 days \hfill needs stop dates};
\node[hit,  below=of a2] (a3)
    {(iii)~~Three or more separate courses \hfill computable: request count};
\node[miss, below=of a3] (a4)
    {(iv)~~Single course $>$ 30 days \hfill needs stop dates};

% legend
\node[draw, very thick, fill=black!12, minimum width=5mm, minimum height=2.5mm,
      inner sep=0pt, anchor=west] (k1) at ([yshift=-7mm] a4.west) {};
\node[font=\scriptsize, right=1mm of k1] (k1t) {computable from the FHIR graph};
\node[draw, dashed, minimum width=5mm, minimum height=2.5mm,
      inner sep=0pt, right=6mm of k1t] (k2) {};
\node[font=\scriptsize, right=1mm of k2] {requires the CSV export};
\end{tikzpicture}
\caption{The four arms of the chronic-NSAID rule. Only arms (i) and (iii) are computable
from the FHIR graph; the duration-based arms require stop dates that the FHIR export does
not carry.}
\label{fig:chronicity}
\end{figure}

\section{Evaluation Method}
\label{sec:eval-method}

The evaluation combines an intrinsic assessment of the graph artifact with an extrinsic
assessment of the recommendation output.

\subsection{Reference standard: a silver standard validated against gold}
\label{sec:oracle}

The evaluation uses two reference sets, and the distinction between them is the one drawn
in the biomedical text-mining literature. A \emph{gold standard} is a reference that has
been \emph{manually} curated; a \emph{silver standard} is one that has been generated
\emph{automatically}. The terms were introduced together by Rebholz-Schuhmann
\textit{et al.} \cite{rebholz2010,rebholz2011}, who contrast a ``manually annotated gold
standard corpus'' with a corpus produced automatically, and who motivate the latter on the
grounds that preparing the former is ``time-consuming and costly'' and therefore does not
scale beyond a few thousand items.

That motivation applies directly here. Labelling every patient in the gate-evaluable
cohort by hand would require a human to adjudicate 876 patients against
81 recommendations---on the order of seventy thousand individual judgements,
which is not feasible for a single investigator. The evaluation therefore adopts the two-tier design used in electronic-health-record
phenotyping. In that design, automatically derived silver labels are produced at scale and
then validated against a small, manually curated gold set \cite{nogues2022}. Here, a silver
standard labels the whole cohort, as described in the remainder of this section. A manually
labelled sample then serves as the gold standard against which that silver standard is
checked (Section~\ref{sec:goldlabel}).

One qualification on the terminology. In the corpus work that introduced it, a silver
standard is built by \emph{harmonising the outputs of several independent automatic
systems}, so that agreement between systems substitutes for human judgement. The silver
standard used here is not a consensus artifact of that kind: it is a single, deterministic
rule-based encoding of the guideline. The term is therefore used in its broader and now
common sense---an automatically generated reference standard---which is how it is applied
in the clinical phenotyping literature \cite{nogues2022}, rather than in the specific
harmonisation sense of the original CALBC corpora.
The term \emph{oracle}, used throughout this thesis, refers to the implementation of that
silver standard. The oracle is the component that computes a label for every
patient--recommendation pair by applying the guideline logic to the raw record, along a path
separate from the system under test. Oracle, silver standard and automated reference standard
therefore name the same component.

\subsubsection*{Construction and two tiers of independence}
The silver standard encodes the guideline logic directly and labels
each patient by reading the \emph{raw FHIR bundle}---never the graph or the CRS output.
That independence has to be stated in two tiers, because not every recommendation gets it to
the same degree, and calling it a single property would claim more than holds.

For the 53 recommendations that are not observation-gated, the two sides differ in both logic
and data path. The system decides by graph traversal---the gate and a conjunction over codes
bridged through UMLS concepts---while the oracle decides by substring matching over condition
names read from the bundle. They disagree often, and that disagreement is what makes precision
and recall informative there.

For the 15 observation-gated recommendations, the independence is of the data path only. The
oracle's numeric thresholds and aggregations---worst-ever laboratory values, most-recent blood
pressure---are used verbatim by the observation gates, so the predicates on the two sides are
not merely similar but the same function. This was checked rather than assumed. Both were evaluated on the same feature vectors: the
876 real cohort vectors and 200{,}000 synthetic ones aimed deliberately at every threshold
they test. Not one disagreement was found (Appendix~\ref{app:modules},
\texttt{predicate\_audit.py}).

The consequence should be stated plainly. For those fifteen, agreement between the system and
the reference standard cannot be evidence that the threshold is clinically right, because both
sides apply the same threshold. What it does establish is that the loader carries patient
values from the bundle into the graph unchanged---a claim about the data path, not about the
rule. An independent test of the rule itself requires the judgement of practising clinicians.
The clinician review (Section~\ref{sec:clinician-method}) provides such a test, but only for
three of those fifteen recommendations; the rest is recorded as a limitation in
Section~\ref{sec:limitations}.
Section~\ref{sec:shared-predicate} reports how much of the concordance figure rests on those
fifteen.

One further documented exception is the chronic-NSAID recommendation, whose duration-based chronicity
arms are read from the Synthea CSV export---still raw synthetic data, never the CRS
output---because the FHIR export omits medication stop dates; this asymmetry and its
consequences are detailed in Section~\ref{sec:chronicity}.
The implementing module is named \texttt{silver\_oracle.py} and writes
\texttt{eval/silver\_oracle\_labels.csv}, while the manually curated labels stay in
\texttt{eval/gold\_labels.csv}; the two filenames therefore carry the same distinction the
text draws. Appendix~\ref{app:modules} lists the modules and what each produces.

\subsection{Guideline-concordance evaluation and the observation-gate ablation}
The primary metric is guideline concordance: the CRS's applicable-recommendation set is
compared with the oracle's labels over the gate-evaluable cohort, reported as precision,
recall, and $F_1$. The central experiment is an \emph{ablation} that runs the system with
and without observation gating; the difference isolates the contribution of observation
gating itself, not that of the container (Section~\ref{sec:csv-gating}). Each recommendation is additionally classified by how
discriminating it is (whether its targets are encoded and bridge to the EHR), which
explains the sources of residual false positives.

\subsection{Manual gold labelling and silver--gold agreement}
\label{sec:goldlabel}
To validate the silver standard, a sample of 30 patients is labelled manually
against the guideline text and the full patient record, producing the \emph{gold standard}.
The silver standard is then compared against this gold set using Cohen's $\kappa$
\cite{cohen1960}, a chance-corrected agreement statistic that has also been used to
benchmark language-model outputs against reference labels in clinical settings
\cite{shah2024}. The uncertainty of $\kappa$ is estimated with a bootstrap \cite{efron1994}
that resamples \emph{patients}, not cells. Cells from one patient are not independent, because
the same record determines all of its labels, and resampling by cluster preserves that
dependence \cite{field2007}. Because $\kappa$ is also affected by label prevalence and by bias
between raters, the prevalence index, the bias index and PABAK are reported alongside it
\cite{byrt1993}. High silver--gold agreement is what licenses the use of the silver
standard across the full cohort. It establishes that the automatically generated labels
reproduce human judgement on the sample where both are available. The concordance
figures computed over all 876 patients therefore rest on a validated reference rather
than an assumed one.
The labelling was performed by the author, reading the guideline text against each
patient's full record. Four decisions recur often enough to have been fixed in advance.
An anticoagulation indication is atrial fibrillation, a mechanical or prosthetic valve,
venous thromboembolism, or an existing warfarin prescription. Dual antiplatelet therapy
is aspirin together with clopidogrel, ticagrelor or prasugrel. An acute coronary syndrome
is STEMI, NSTEMI or preinfarction syndrome. Statin intensity follows the guideline's own bands. Atorvastatin~40--80~mg and
rosuvastatin~20~mg count as high-intensity. Atorvastatin~10--20~mg, rosuvastatin~5--10~mg,
simvastatin~20--40~mg, pravastatin~40--80~mg and ezetimibe/simvastatin~40~mg count as
moderate. Simvastatin~10~mg, pravastatin~20~mg and lovastatin~10~mg count as low. The number of recommendations labelled is not uniform across patients: it
ranges from 71 to 73, for 2{,}160 cells in total. Three contraindication recommendations
(\texttt{APT:7}, \texttt{APT:9} and \texttt{APT:10}) are labelled only for patients
actually on the drug concerned, since for anyone else the recommendation has no subject to
judge. A single annotator
means no inter-rater reliability can be reported for the gold standard itself, which is
revisited among the limitations in Section~\ref{sec:limitations}.

\subsection{Safety-detection evaluation}
Contraindication-safety detection is evaluated on the manually labelled safety flags. The
unary and population-gated variants are each scored for precision and recall. Together the
two scores illustrate the trade-off between over-flagging in the unary variant and missing
structurally un-gateable contraindications in the population-gated one. The chronic-NSAID recommendation is scored separately
under its chronicity-gated variant, which quantifies the recall cost of relying on the FHIR
categorical medication status once the duration-based chronicity arms are unavailable
(Section~\ref{sec:chronicity}).

\subsection{Briefing faithfulness evaluation}
The grounded briefing is evaluated for faithfulness by an automated procedure that compares
the generated text against the retrieved facts. The procedure flags any fabricated
recommendation identifier or PMID, any omitted applicable recommendation or safety flag, and
one over-claim pattern: stating that the patient requires oral anticoagulation when no
anticoagulant is on the medication list. A briefing passes when all four counts are zero. The
procedure does not check whether every other sentence is supported by the facts, since that
requires a human reading.

\subsection{Clinician review of recommendation appropriateness}
\label{sec:clinician-method}
The tiers above test the system against rules derived from the guideline. The clinician
review tests what those rules cannot: whether the recommendations the system outputs are,
in a clinician's judgement, appropriate for the patient in question. What is judged is the
system's output, not the gold-standard labels.

\subsubsection*{Review sheet}
The review sheet covers 20 patients, the 20 highest-scoring of the 30 gold-standard sample
patients under that sample's selection rule (Section~\ref{sec:goldlabel}). Each patient
occupies one page with a short identification, the conclusion of the grounded briefing, and
a table of recommendations. The conclusion is the only basis for judgement. The clinical
vignette and the full medication list are deliberately not shown, so that raters judge what
the system actually conveys rather than the record behind it.

The conclusions were translated into Indonesian. The translation was checked mechanically:
citations, blood-pressure targets, body-mass-index thresholds, stages, recommendation classes,
and drug names had to carry over. The author then edited 12 of the 20 translations. The edits
removed bracketed citations and class markers, removed the phrase ``yang sudah tegak''
(``established''), restored a few terms to their English form, and rephrased a few phrases.
Raters therefore judged edited text, not the model's output as generated.

The table lists every safety flag for the patient and every applicable recommendation
triggered by the patient's own conditions. Standard CCD care that applies equally to every
patient is not listed. In all there are 327 rows covering 21 recommendation types in three
groups: check/stop of a current medication (107 rows), non-drug management (77 rows), and
drug therapy (143 rows). Recommendation codes and classes are not shown; both are kept in a
key file that was not sent to raters. Each row is judged with one of three options:
``Tepat'' (appropriate), ``Belum tepat'' (not yet appropriate), or ``Tidak dapat dinilai''
(cannot be judged). The rating guide states that the third option is a valid answer, because
the conclusion does not always carry all the information needed. Raters were asked to write
a comment on every row they did not judge appropriate.

\subsubsection*{Raters and coverage}
Three clinicians rated independently and were asked not to confer before finishing. Raters~A and~C are
cardiologists (\emph{spesialis jantung dan pembuluh darah}, SpJP), C the more senior. Rater~B
is a general practitioner enrolled in a doctoral programme. None completed the whole sheet.
Rater~A rated patients 1--15, rater~B patients 1--10, and rater~C patient~1 only. Within
that coverage, raters~A and~B each skipped one row. The rows rated are therefore 247 for
rater~A, 160 for rater~B, and 16 for rater~C.

\subsubsection*{Analysis}
The proportion of each option is computed per rater. A 95\% confidence interval for the
proportion rated appropriate is estimated by a 10{,}000-replicate bootstrap that resamples
patients, for the same reason as for the silver-gold $\kappa$
(Section~\ref{sec:goldlabel}). Agreement between raters~A and~B is measured on the rows both
rated, as observed agreement, three-category Cohen's $\kappa$ \cite{cohen1960}, and PABAK
\cite{byrt1993}, with per-patient bootstrap intervals. Rater~C covers a single patient and is
reported descriptively. Comments are read and grouped by recommendation. All figures are
computed by \texttt{baca\_hasil.py} and \texttt{analisis\_hasil.py}
(Appendix~\ref{app:modules}).

\subsection{Graph structural evaluation}
\label{sec:structural}
The evaluation so far is \emph{extrinsic}: it judges the graph by the quality of the
decisions taken over it. The graph is also assessed \emph{intrinsically}, as an artifact in
its own right, because a recommendation output says nothing about whether the underlying
representation is heterogeneous, broad, or deeply connected enough to be reused.

This assessment adopts the framing established for ontology evaluation by Gangemi \textit{et
al.} \cite{gangemi2006}. They distinguish three kinds of measure. \emph{Structural} measures
treat the ontology as a graph. Functional measures concern its fitness for an intended use,
and usability-profiling measures concern its annotation. The tier described here is
structural in that sense, whereas the concordance, safety, faithfulness, and clinician-review
tiers above are functional. Structural metrics of this family have been evaluated empirically over
public ontology repositories and none was found to produce an unstable clustering of them: across 119 OBO~Foundry
and 78 AgroPortal ontologies, the classification each metric induces survives bootstrap
resampling of the corpus, which is the sense in which these metrics are stable
\cite{franco2020}.
The same argument this thesis makes---that scale counts alone do not
characterise a knowledge graph---motivates the structural metric sets proposed for
cross-domain knowledge graphs \cite{seo2022}.

Three dimensions are reported. \textbf{Richness} captures schema and vocabulary
heterogeneity: the number of node labels, relationship types, and code systems, and the
number of fused semantic layers. This follows the notion of richness introduced by OntoQA
\cite{tartir2005}, which characterises an ontology through the variety of relationships and
classes actually used rather than through raw size. OntoQA states that notion as a
ratio---its relationship richness is $P/(SC+P)$, the share of non-inheritance relations
among all connections between classes---which this graph cannot form, having no subsumption
hierarchy and therefore no $SC$ to divide by. What is reported here is the variety itself,
as counts. \textbf{Width} captures breadth of
coverage: node and edge cardinality, cohort size, average conditions per patient, and the
proportion of the guideline that is represented. Width in this sense corresponds to the
\emph{completeness} dimension of knowledge-graph quality, for which a substantial
methodological literature exists \cite{issa2021,xue2023}. \textbf{Depth} captures how far a
conclusion can be traced: the length of the canonical
patient-to-recommendation-to-evidence inference path, and the degree to which shared
\texttt{:Code} nodes act as hubs.

Two departures from the source literature must be stated, because the terms are shared but
the quantities are not. First, \emph{depth} in the ontology-evaluation literature normally
means the depth of the subsumption hierarchy---how many \texttt{is-a} levels a class tree
has. The graph built here is a property graph with no subsumption hierarchy to measure, so
depth is redefined as the length of an inference path across the fused layers. Second,
ontology \emph{breadth} normally means the number of sibling classes at a level, whereas
width here is cardinality and cohort coverage. What is borrowed is therefore the
\emph{classification} of measures rather than the formulas: Gangemi's separation of structural
from functional evaluation \cite{gangemi2006}, and OntoQA's principle of characterising a graph
by the variety of its relationships rather than by its size \cite{tartir2005}. The three
quantities reported in Chapter~IV are defined by this thesis, and should not be compared
numerically against values published elsewhere under the same names.

Cohen's $\kappa$ and the concordance metrics are computed directly from the
$2\times2$ agreement table by the evaluation scripts themselves, without a statistics
library, so reproducing them requires no dependency beyond the Python standard library.

\section{Research Ethics, Problems, and Limitations}
\label{sec:ethics}

\subsection{Ethical considerations}
The patient data used in this research are entirely synthetic. Synthea generates
statistically plausible but fictitious individuals \cite{synthea}, so no protected health
information is processed, no patient consent is required, and there is no risk of
re-identification. This is a deliberate design choice rather than a convenience. It allows the full record,
including medication histories and laboratory trajectories, to be published alongside the
code for reproduction. That would be impossible with real records, which are mostly still
unstructured narrative text.

One component of this research involves people: the clinician review
(Section~\ref{sec:clinician-method}). The clinicians act as expert raters of the system's
output, not as research subjects. What is collected is their judgement of recommendations for
synthetic patients, not data about themselves. Participation was voluntary, and raters are
identified only by code in the thesis and in any released data. Ethical clearance for
human-subject research is therefore not required. The remaining evaluation tiers are
automated: the reference oracle is generated programmatically, the gold standard is labelled
by the author from synthetic records, and briefing faithfulness is checked by an automated
verification procedure. The consequence for the strength of the evidence is stated
explicitly as a limitation in Section~\ref{sec:limitations}.

A third consideration is the intended role of the system. It is designed as a
decision-\emph{support} tool that surfaces guideline recommendations and contraindication
flags for a clinician to judge; it does not prescribe, and it is not validated for
autonomous use. The grounded briefing is explicitly constrained to phrase retrieved facts
rather than to reason clinically, and every statement it makes carries its recommendation
identifier and class of recommendation so that a reader can audit it against the source
guideline.

\subsection{Problems and limitations}
\label{sec:limitations}
The following limitations are inherent to the design and constrain how the results should
be read:

\begin{enumerate}
    \item \textbf{Synthetic data.} Synthea patients are generated from disease-progression
    models, so their comorbidity patterns and laboratory trajectories are internally
    consistent but not drawn from a real population. Performance measured here therefore
    establishes that the pipeline is \emph{internally} correct, not that it would achieve
    the same concordance on real records, which are noisier, sparser, and contain
    unstructured text.
    \item \textbf{No cardiac symptom data.} Synthea records carry some patient-reported
    measures: a pain score, and questionnaires on depression, anxiety, substance use and social
    needs. They carry no description of cardiac symptoms. Across the 30 sample patients, each
    patient's chief-complaint list never changes over 8{,}067 clinical notes. Only one patient
    lists a cardiac symptom, chest pressure and shortness of breath, and there is no angina
    grade, angina frequency or exercise capacity anywhere. An NYHA functional class is recorded
    once for five patients, but it lies outside the observation allow-list. None of the 876
    cohort patients carries a coded angina diagnosis. As a result, the six recommendations that
    depend on angina (\texttt{CR:1:2}, \texttt{MTRA:1}--\texttt{MTRA:4} and \texttt{MRA:1})
    are never positive, so this part of the guideline goes entirely untested. In real records
    such symptoms are usually recorded as free text, and capturing them requires structured
    encoding that this design does not cover.
    \item \textbf{Information loss in the FHIR projection.} As set out in
    Section~\ref{sec:chronicity}, the FHIR export carries a categorical medication status
    but no stop dates, so duration-dependent guideline logic cannot be evaluated from the
    graph alone. This is a property of the export, not of the implementation, and it is
    quantified rather than hidden.
    \item \textbf{Shared predicate on the observation-gated recommendations.} For 15 of the 68
    recommendations, the system's predicate and the reference standard's rule are the same
    function (Section~\ref{sec:oracle}). Perfect agreement in the two gated classes therefore
    tests the data path, not the correctness of the rule. The part scored by a genuinely
    independent judge covers 53 recommendations and reaches precision $0.666$ and $F_1$
    $0.797$, against $0.685$ and $0.811$ for the headline
    (Section~\ref{sec:shared-predicate}). The whole of the ablation gain is measured on those
    fifteen shared-predicate recommendations, so what it measures is implementation fidelity
    rather than clinical validity. The clinician review tests the rule for only three of the
    fifteen, \texttt{WM:2}, \texttt{WM:3}, and \texttt{SGLT2:4}; all 50 judgements of them by
    raters~A and~B are ``appropriate'' (Section~\ref{sec:clinician}).
    \item \textbf{Single gold-standard annotator.} All 2{,}160 gold-standard cells were
    labelled by one person, the author, who also encoded the silver standard's rules. No
    inter-rater reliability can therefore be reported for the gold standard. High agreement may
    also partly reflect a shared reading of the guideline text, not only the correctness of the
    rules. The confidence interval in Section~\ref{sec:kappa-ci} addresses the uncertainty due
    to sample size but not this limitation, since resampling cannot produce a second rater.
    Relabelling by a second annotator independent of the author is the only way to close it.
    \item \textbf{FHIR is not a precondition for the observation gate.} The precision gain
    reported in Section~\ref{sec:ablation} comes from gating on numeric observation values, not
    from the FHIR container specifically. The same gate can be rebuilt from Synthea's tabular
    export and returns identical decisions on all 13{,}140 decisions compared
    (Section~\ref{sec:csv-gating}). The distinction bounds what the FHIR-native design may
    claim: directness and convenience, not feasibility. It should be added that this result
    holds for Synthea's export, which carries an \texttt{ENCOUNTER} foreign key on every
    observation row; a tabular export without one need not fare the same way.
    \item \textbf{Single guideline, single disease.} Only the 2023 AHA/ACC CCD guideline is
    encoded, and only for chronic coronary disease. Generalisation to other guidelines or
    conditions is plausible given the concept-bridge design but is not demonstrated.
    \item \textbf{Single-concept cohort gate.} As set out in Section~\ref{sec:gate}, the CCD
    gate accepts one UMLS concept, whereas the guideline defines chronic coronary disease
    disjunctively. Sixty-seven patients ($7.1\%$ of the clinically defined cohort) whose only
    CCD evidence is a prior-infarction or prior-revascularization \emph{situation} code are
    consequently never matched against any recommendation. This bounds the reported results
    to a conservative subset and would, in a deployment, cause those patients to receive no
    output at all. Widening the gate to a disjunction of concepts is a small change to the
    query and is the first item of future work.
    \item \textbf{No prospective or outcome validation.} Concordance is measured against an
    automatically generated silver standard and a manually curated gold standard, not
    against patient outcomes. Nothing here establishes that following the system's
    recommendations improves care.
    \item \textbf{Stochastic phrasing.} The briefing's wording is produced by a language
    model and varies between runs, so faithfulness is verified per run rather than proved
    once. Retrieval, by contrast, is deterministic and reproducible.
    \item \textbf{Scope of the faithfulness check.} The automated check covers 30 briefings
    and tests citable items, recommendation completeness, and one over-claim pattern. The
    semantic faithfulness of every sentence is not tested. What is supported, therefore, is
    that the briefings neither fabricate recommendations or citations nor omit
    recommendations, not that every sentence is faithful to the facts. The over-claim test
    itself was corrected after its output was seen (Section~\ref{sec:faithfulness}). Thirty
    briefings are also still too few to estimate an error rate precisely.
    \item \textbf{A limited clinician review.} The reference oracle, the safety-detection
    check, and the faithfulness verification test the system against \emph{rules derived
    from the guideline}. Only the clinician review (Section~\ref{sec:clinician}) tests it
    against clinical judgement, and its coverage is limited. Rater~A, the only cardiologist
    to rate more than one patient, covered 15 of the 20 patients. Rater~B, a general
    practitioner, covered 10, and rater~C a single patient. Agreement between raters~A and~B
    is fair ($\kappa = 0.29$), so the proportion judged appropriate depends on who rates.
    Raters also read only the conclusion and the recommendation table, not the full record.
    Some ``not yet appropriate'' and ``cannot be judged'' ratings name information that was
    not shown, such as ejection fraction, heart-failure status, bleeding risk, and the
    medication list. The conclusion the raters read was, moreover, a translation edited by
    the author. The review therefore locates the system's weak points but does not yet
    estimate, in general, the gap between guideline concordance and clinical
    appropriateness. It also judges the appropriateness of recommendations, not the quality
    of the summary. A reader study on the model of Van Veen \textit{et al.}~\cite{vanveen2023},
    in which clinicians score summaries for completeness, correctness, and conciseness, with
    inter-rater agreement quantified by Cohen's $\kappa$ \cite{cohen1960}, is still left to
    future work.
\end{enumerate}
